% Options for packages loaded elsewhere
\PassOptionsToPackage{unicode}{hyperref}
\PassOptionsToPackage{hyphens}{url}
%
\documentclass[
]{article}
\usepackage{amsmath,amssymb}
\usepackage{iftex}
\ifPDFTeX
  \usepackage[T1]{fontenc}
  \usepackage[utf8]{inputenc}
  \usepackage{textcomp} % provide euro and other symbols
\else % if luatex or xetex
  \usepackage{unicode-math} % this also loads fontspec
  \defaultfontfeatures{Scale=MatchLowercase}
  \defaultfontfeatures[\rmfamily]{Ligatures=TeX,Scale=1}
\fi
\usepackage{lmodern}
\ifPDFTeX\else
  % xetex/luatex font selection
\fi
% Use upquote if available, for straight quotes in verbatim environments
\IfFileExists{upquote.sty}{\usepackage{upquote}}{}
\IfFileExists{microtype.sty}{% use microtype if available
  \usepackage[]{microtype}
  \UseMicrotypeSet[protrusion]{basicmath} % disable protrusion for tt fonts
}{}
\makeatletter
\@ifundefined{KOMAClassName}{% if non-KOMA class
  \IfFileExists{parskip.sty}{%
    \usepackage{parskip}
  }{% else
    \setlength{\parindent}{0pt}
    \setlength{\parskip}{6pt plus 2pt minus 1pt}}
}{% if KOMA class
  \KOMAoptions{parskip=half}}
\makeatother
\usepackage{xcolor}
\usepackage[margin=1in]{geometry}
\usepackage{listings}
\newcommand{\passthrough}[1]{#1}
\lstset{defaultdialect=[5.3]Lua}
\lstset{defaultdialect=[x86masm]Assembler}
\usepackage{longtable,booktabs,array}
\usepackage{calc} % for calculating minipage widths
% Correct order of tables after \paragraph or \subparagraph
\usepackage{etoolbox}
\makeatletter
\patchcmd\longtable{\par}{\if@noskipsec\mbox{}\fi\par}{}{}
\makeatother
% Allow footnotes in longtable head/foot
\IfFileExists{footnotehyper.sty}{\usepackage{footnotehyper}}{\usepackage{footnote}}
\makesavenoteenv{longtable}
\setlength{\emergencystretch}{3em} % prevent overfull lines
\providecommand{\tightlist}{%
  \setlength{\itemsep}{0pt}\setlength{\parskip}{0pt}}
\setcounter{secnumdepth}{5}
\ifLuaTeX
  \usepackage{selnolig}  % disable illegal ligatures
\fi
\usepackage{bookmark}
\IfFileExists{xurl.sty}{\usepackage{xurl}}{} % add URL line breaks if available
\urlstyle{same}
\hypersetup{
  hidelinks,
  pdfcreator={LaTeX via pandoc}}

\author{}
\date{}

\begin{document}

{
\setcounter{tocdepth}{3}
\tableofcontents
}
\section{Destructuring Field Guide}\label{destructuring-field-guide}

Name: \textbf{Zeynep Sudem Aktaş}\\
Student number: \textbf{2025502009}

Bonuses claimed: \textbf{Bonus 1: LaTeX; Bonus 2: Depth; Bonus 3: Real
code}. Bonus 4 (Git repository) is not claimed yet.

\subsection{Run}\label{run}

\begin{lstlisting}[language=bash]
npm install
npm run all      # runs the JavaScript and TypeScript examples
npm run check    # type-checks everything in strict mode
npm run all > output.txt
\end{lstlisting}

Use Node.js 20.19 or newer. Do not submit
\passthrough{\lstinline!node\_modules!}.

\subsection{What destructuring does}\label{what-destructuring-does}

Destructuring takes values out of an array or an object and assigns them
to variables with one pattern on the left side of an assignment. It is
shorter than writing one property or index lookup for every variable.
Array patterns use positions; object patterns use property names.

A useful way to understand any pattern is to rewrite it as ordinary
assignments. For example, \passthrough{\lstinline!const [x, y] = pair!}
is conceptually similar to
\passthrough{\lstinline!const x = pair[0]; const y = pair[1]!}.
Likewise, \passthrough{\lstinline!const \{ title \} = task!} is similar
to \passthrough{\lstinline!const title = task.title!}.

\subsection{Field guide}\label{field-guide}

\subsubsection{\texorpdfstring{Form 1: array values by position,
\texttt{const\ {[}x,\ y{]}\ =\ f()}}{Form 1: array values by position, const {[}x, y{]} = f()}}\label{form-1-array-values-by-position-const-x-y--f}

\textbf{Syntax:}
\passthrough{\lstinline!const [first, second] = someArray!}.

\textbf{Example} (\passthrough{\lstinline!js/arrays.mjs!}):

\begin{lstlisting}
function minMax(numbers) {
  return [Math.min(...numbers), Math.max(...numbers)]
}
const [lowest, highest] = minMax([4, 9, 2])
console.log('form 1:', lowest, highest)
\end{lstlisting}

\textbf{Output:} \passthrough{\lstinline!form 1: 2 9!}

\textbf{Use it when:} a function returns a small group of values whose
positions have a clear meaning, such as a minimum and maximum.

\textbf{Watch out:} positions matter. If the returned array order
changes, the variables receive different values. A missing position
becomes \passthrough{\lstinline!undefined!} unless a default is
supplied.

\subsubsection{\texorpdfstring{Form 2: skip an item and collect the
rest,
\texttt{const\ {[},\ second,\ ...others{]}\ =\ list}}{Form 2: skip an item and collect the rest, const {[}, second, ...others{]} = list}}\label{form-2-skip-an-item-and-collect-the-rest-const--second-others--list}

\textbf{Syntax:} commas leave positions unused;
\passthrough{\lstinline!...name!} in a pattern collects the remaining
elements into a new array.

\textbf{Example} (\passthrough{\lstinline!js/arrays.mjs!}):

\begin{lstlisting}
const readingList = ['The Hobbit', 'Dune', 'Klara and the Sun', 'Piranesi']
const [, currentBook, ...laterBooks] = readingList
console.log('form 2:', currentBook, laterBooks)
\end{lstlisting}

\textbf{Output:}
\passthrough{\lstinline!form 2: Dune [ 'Klara and the Sun', 'Piranesi' ]!}

\textbf{Use it when:} the first value is irrelevant, but later positions
or all remaining values are useful.

\textbf{Watch out:} the rest element must be last in its pattern. It
collects the remaining items; it does not include the skipped first
item. Rest in a pattern collects, while spread in an expression expands.

\subsubsection{\texorpdfstring{Form 3: swap two variables,
\texttt{{[}a,\ b{]}\ =\ {[}b,\ a{]}}}{Form 3: swap two variables, {[}a, b{]} = {[}b, a{]}}}\label{form-3-swap-two-variables-a-b--b-a}

\textbf{Syntax:}
\passthrough{\lstinline![left, right] = [right, left]!}.

\textbf{Example} (\passthrough{\lstinline!js/arrays.mjs!}):

\begin{lstlisting}
let firstStop = 'Oslo'
let secondStop = 'Copenhagen'
;[firstStop, secondStop] = [secondStop, firstStop]
console.log('form 3:', firstStop, secondStop)
\end{lstlisting}

\textbf{Output:} \passthrough{\lstinline!form 3: Copenhagen Oslo!}

\textbf{Use it when:} two already-declared variables need to exchange
values.

\textbf{Watch out:} this is assignment destructuring, not a declaration,
so the variables must already exist. Parentheses are needed when an
object pattern starts an assignment statement; the leading semicolon
here protects against automatic semicolon insertion after a preceding
line.

\subsubsection{\texorpdfstring{Form 4: array defaults,
\texttt{const\ {[}a\ =\ 1,\ b\ =\ 2{]}\ =\ values}}{Form 4: array defaults, const {[}a = 1, b = 2{]} = values}}\label{form-4-array-defaults-const-a--1-b--2--values}

\textbf{Syntax:} put \passthrough{\lstinline!= defaultValue!} after a
pattern position.

\textbf{Example} (\passthrough{\lstinline!js/arrays.mjs!}):

\begin{lstlisting}
const [priority = 'normal', reminder = 'tomorrow'] = ['high']
console.log('form 4:', priority, reminder)
\end{lstlisting}

\textbf{Output:} \passthrough{\lstinline!form 4: high tomorrow!}

\textbf{Use it when:} an array may omit a position or contain
\passthrough{\lstinline!undefined!}, and a useful fallback is needed.

\textbf{Watch out:} the default is used for
\passthrough{\lstinline!undefined!}, including a missing array position.
It is not used for \passthrough{\lstinline!null!},
\passthrough{\lstinline!false!}, \passthrough{\lstinline!0!}, or an
empty string.

\subsubsection{\texorpdfstring{Form 5: object properties by name,
\texttt{const\ \{\ a,\ b\ \}\ =\ object}}{Form 5: object properties by name, const \{ a, b \} = object}}\label{form-5-object-properties-by-name-const--a-b---object}

\textbf{Syntax:} property names on the left match keys on the object.

\textbf{Example} (\passthrough{\lstinline!js/objects.mjs!}):

\begin{lstlisting}
const task = { id: 't1', title: 'Buy milk', status: 'todo' }
const { title, status } = task
console.log('form 5:', title, status)
\end{lstlisting}

\textbf{Output:} \passthrough{\lstinline!form 5: Buy milk todo!}

\textbf{Use it when:} several known properties of an object are needed
in the current scope.

\textbf{Watch out:} unlike array destructuring, property order does not
matter. The property names must match the object keys unless a rename is
used. The source must not be \passthrough{\lstinline!null!} or
\passthrough{\lstinline!undefined!}.

\subsubsection{\texorpdfstring{Form 6: rename a property,
\texttt{const\ \{\ a:\ alpha\ \}\ =\ object}}{Form 6: rename a property, const \{ a: alpha \} = object}}\label{form-6-rename-a-property-const--a-alpha---object}

\textbf{Syntax:}
\passthrough{\lstinline!\{ originalProperty: newVariableName \}!}.

\textbf{Example} (\passthrough{\lstinline!js/objects.mjs!}):

\begin{lstlisting}
const book = { title: 'Dune', year: 1965 }
const { title: bookName } = book
console.log('form 6:', bookName)
\end{lstlisting}

\textbf{Output:} \passthrough{\lstinline!form 6: Dune!}

\textbf{Use it when:} the property name is unclear in the current
context or conflicts with an existing variable name.

\textbf{Watch out:} \passthrough{\lstinline!title!} is the source
property and \passthrough{\lstinline!bookName!} is the variable created.
There is no new variable called \passthrough{\lstinline!title!} in this
example. Read the colon as ``take \passthrough{\lstinline!title!} and
call it \passthrough{\lstinline!bookName!}.''

\subsubsection{\texorpdfstring{Form 7: object defaults,
\texttt{const\ \{\ a\ =\ 1\ \}\ =\ object}}{Form 7: object defaults, const \{ a = 1 \} = object}}\label{form-7-object-defaults-const--a--1---object}

\textbf{Syntax:} add a default after the property variable.

\textbf{Example} (\passthrough{\lstinline!js/objects.mjs!}):

\begin{lstlisting}
const event = { name: 'MIS study session' }
const { location = 'Library', duration = 60 } = event
console.log('form 7:', location, duration)
\end{lstlisting}

\textbf{Output:} \passthrough{\lstinline!form 7: Library 60!}

\textbf{Use it when:} an object property is optional and the code needs
a fallback value.

\textbf{Watch out:} the default applies when the property is missing or
its value is \passthrough{\lstinline!undefined!}. If the property is
explicitly \passthrough{\lstinline!null!}, the result stays
\passthrough{\lstinline!null!}.

\subsubsection{\texorpdfstring{Form 8: object rest,
\texttt{const\ \{\ a,\ ...others\ \}\ =\ object}}{Form 8: object rest, const \{ a, ...others \} = object}}\label{form-8-object-rest-const--a-others---object}

\textbf{Syntax:} extract selected properties first;
\passthrough{\lstinline!...others!} collects the remaining enumerable
own properties into a new object.

\textbf{Example} (\passthrough{\lstinline!js/objects.mjs!}):

\begin{lstlisting}
const profile = { id: 'u7', username: 'zey', department: 'MIS', year: 2 }
const { id, ...profileDetails } = profile
console.log('form 8:', id, profileDetails)
\end{lstlisting}

\textbf{Output:}
\passthrough{\lstinline!form 8: u7 \{ username: 'zey', department: 'MIS', year: 2 \}!}

\textbf{Use it when:} one or more properties should be separated from
the remaining object, for example before forwarding the remaining props.

\textbf{Watch out:} the rest object is a shallow copy of the remaining
properties. Nested objects inside it are still shared references. Rest
must be the last item in the object pattern.

\subsubsection{\texorpdfstring{Form 9: nested patterns,
\texttt{\{\ owner:\ \{\ name\ \},\ milestones:\ {[}first{]}\ \}}}{Form 9: nested patterns, \{ owner: \{ name \}, milestones: {[}first{]} \}}}\label{form-9-nested-patterns--owner--name--milestones-first-}

\textbf{Syntax:} put another object or array pattern inside a property
pattern.

\textbf{Example} (\passthrough{\lstinline!js/objects.mjs!}):

\begin{lstlisting}
const project = {
  owner: { name: 'Ada', role: 'lead' },
  milestones: ['prototype', 'testing', 'launch']
}
const { owner: { name: ownerName }, milestones: [firstMilestone, ...remainingMilestones] } = project
console.log('form 9:', ownerName, firstMilestone, remainingMilestones)
\end{lstlisting}

\textbf{Output:}
\passthrough{\lstinline!form 9: Ada prototype [ 'testing', 'launch' ]!}

\textbf{Use it when:} the required values are nested inside an object or
an array stored in an object property.

\textbf{Watch out:} every intermediate object or array must exist. If
\passthrough{\lstinline!owner!} is \passthrough{\lstinline!undefined!},
trying to destructure \passthrough{\lstinline!name!} from it throws a
\passthrough{\lstinline!TypeError!}. Deep patterns can become difficult
to read; split them into steps when necessary.

\subsubsection{\texorpdfstring{Form 10: computed property key,
\texttt{const\ \{\ {[}key{]}:\ value\ \}\ =\ object}}{Form 10: computed property key, const \{ {[}key{]}: value \} = object}}\label{form-10-computed-property-key-const--key-value---object}

\textbf{Syntax:} put an expression in square brackets to select the
property dynamically.

\textbf{Example} (\passthrough{\lstinline!js/objects.mjs!}):

\begin{lstlisting}
const selectedField = 'status'
const ticket = { id: 't9', status: 'in-progress', title: 'Fix login' }
const { [selectedField]: selectedValue } = ticket
console.log('form 10:', selectedField, selectedValue)
\end{lstlisting}

\textbf{Output:} \passthrough{\lstinline!form 10: status in-progress!}

\textbf{Use it when:} the property to extract is stored in a variable or
computed at runtime.

\textbf{Watch out:} the brackets evaluate the key expression. The colon
still renames the selected property into a local variable; it does not
create a variable named after the expression.

\subsubsection{\texorpdfstring{Form 11: destructured function
parameters,
\texttt{function\ f(\{\ a\ =\ 1\ \}\ =\ \{\})}}{Form 11: destructured function parameters, function f(\{ a = 1 \} = \{\})}}\label{form-11-destructured-function-parameters-function-f-a--1---}

\textbf{Syntax:} destructure the options object in the parameter list.
The type annotation in TypeScript belongs after the complete pattern.

\textbf{Example} (\passthrough{\lstinline!js/params.mjs!}):

\begin{lstlisting}
function createTask({ title, status = 'todo' } = {}) {
  return { title, status }
}
console.log('form 11:', createTask({ title: 'Read' }))
\end{lstlisting}

\textbf{Output:}
\passthrough{\lstinline!form 11: \{ title: 'Read', status: 'todo' \}!}

\textbf{Use it when:} a function accepts named options, especially when
callers may omit some options.

\textbf{Watch out:} the \passthrough{\lstinline!= \{\}!} defaults the
whole parameter when the argument is
\passthrough{\lstinline!undefined!}; it does not protect against
\passthrough{\lstinline!null!}. A default on an individual property is a
separate default. The array-parameter variation in the source is Form
11b.

\subsubsection{\texorpdfstring{Form 12: destructuring in loops,
\texttt{for\ (const\ {[}key,\ value{]}\ of\ entries)}}{Form 12: destructuring in loops, for (const {[}key, value{]} of entries)}}\label{form-12-destructuring-in-loops-for-const-key-value-of-entries}

\textbf{Syntax:} each iteration can destructure the pair yielded by
\passthrough{\lstinline!Object.entries(object)!} or a
\passthrough{\lstinline!Map!}.

\textbf{Example} (\passthrough{\lstinline!js/params.mjs!}):

\begin{lstlisting}
const settings = { theme: 'dark', notifications: true }
const settingSummary = []
for (const [key, value] of Object.entries(settings)) {
  settingSummary.push(`${key}=${value}`)
}
console.log('form 12a:', settingSummary.join(', '))
\end{lstlisting}

\textbf{Output:}
\passthrough{\lstinline!form 12a: theme=dark, notifications=true!}

The same pair pattern works with a \passthrough{\lstinline!Map!}:

\begin{lstlisting}
const seatAssignments = new Map([['A1', 'Mina'], ['A2', 'Eren']])
for (const [seat, student] of seatAssignments) {
  // seat is the key; student is the value
}
\end{lstlisting}

\textbf{Output:} \passthrough{\lstinline!form 12b: A1:Mina, A2:Eren!}

\textbf{Use it when:} iterating over key-value pairs and both parts are
needed in the loop body.

\textbf{Watch out:} \passthrough{\lstinline!Object.entries!} returns
string-keyed pairs for an object\textquotesingle s own enumerable string
properties. A \passthrough{\lstinline!Map!} iterator yields its
key-value pairs directly. Do not confuse
\passthrough{\lstinline!for...of!} (iterates values) with
\passthrough{\lstinline!for...in!} (iterates enumerable property keys).

\subsection{Pitfalls}\label{pitfalls}

These examples are in \passthrough{\lstinline!js/pitfalls.mjs!}.

\begin{enumerate}
\def\labelenumi{\arabic{enumi}.}
\tightlist
\item
  \textbf{Destructuring \passthrough{\lstinline!undefined!} or
  \passthrough{\lstinline!null!} throws.} Output:
  \passthrough{\lstinline!pitfall 1: TypeError!}. The right-hand value
  must be an object or array that can be read; check it or provide a
  safe default such as \passthrough{\lstinline!= \{\}!}.
\item
  \textbf{A default does not replace \passthrough{\lstinline!null!}.}
  Output: \passthrough{\lstinline!pitfall 2: none null!}. A
  destructuring default is triggered only by
  \passthrough{\lstinline!undefined!}, not by every falsy value.
\item
  \textbf{Object assignment patterns need parentheses at the start of a
  statement.} Output: \passthrough{\lstinline!pitfall 3: 3!}. Without
  parentheses, \passthrough{\lstinline!\{ count \}!} at the beginning of
  a statement can be parsed as a block instead of an object pattern.
\item
  \textbf{Destructuring is shallow, not a deep clone.} Output:
  \passthrough{\lstinline!pitfall 4: dark!}. The extracted
  \passthrough{\lstinline!settings!} variable refers to the same nested
  object, so mutating it also changes the nested object reached through
  \passthrough{\lstinline!original!}.
\item
  \textbf{Object patterns use names; array patterns use positions.}
  Output:
  \passthrough{\lstinline!pitfall 5: Istanbul Mina | Mina Istanbul!}.
  Reordering object properties in the pattern does not change which key
  is read, but reordering array positions changes the assigned values.
\end{enumerate}

\subsection{Props and children (P1-P7)}\label{props-and-children-p1-p7}

A React function component receives one props object. JSX attributes
become properties on that object. Content between the
component\textquotesingle s opening and closing tags is supplied through
the special property named \passthrough{\lstinline!children!}. It is not
a second argument.

\subsubsection{P1: props is one object}\label{p1-props-is-one-object}

\passthrough{\lstinline!CardA(props)!} reads properties from the object.
\passthrough{\lstinline!CardB(\{ title, children \})!} destructures that
same object. The second form combines Form 5 (object properties by name)
with Form 11 (destructuring in a parameter list).

\textbf{Rendered output} (run \passthrough{\lstinline!npm run all!} to
verify):

\begin{lstlisting}
P1 A: <section><h2>Reading</h2>Book notes</section>
P1 B: <section><h2>Reading</h2>Book notes</section>
\end{lstlisting}

\subsubsection{P2: children is a prop}\label{p2-children-is-a-prop}

\passthrough{\lstinline!<CardB title="Note">Remember the deadline</CardB>!}
and
\passthrough{\lstinline!<CardB title="Note" children="Remember the deadline" />!}
provide the same \passthrough{\lstinline!children!} property in two
different JSX spellings.

\textbf{Rendered output:}

\begin{lstlisting}
P2 nested: <section><h2>Note</h2>Remember the deadline</section>
P2 prop: <section><h2>Note</h2>Remember the deadline</section>
\end{lstlisting}

Do not provide both a \passthrough{\lstinline!children!} attribute and
nested JSX content at once; that creates two competing ways to supply
the same prop and React reports a warning.

\subsubsection{P3: what children can
contain}\label{p3-what-children-can-contain}

React can render strings, numbers, elements, arrays of renderable nodes,
and combinations of them. \passthrough{\lstinline!undefined!},
\passthrough{\lstinline!null!}, and booleans such as
\passthrough{\lstinline!false!} render nothing. A function is not itself
a normal renderable node; in this example it is a \emph{render prop} and
the component deliberately calls it to obtain a React node.

Expected server-rendered outputs from
\passthrough{\lstinline!renderToStaticMarkup!}:

\begin{lstlisting}
P3 string: <div>hello</div>
P3 number: <div>42</div>
P3 one element: <div><strong>bold</strong></div>
P3 several: <div><i>one</i><b>two</b></div>
P3 undefined: <div></div>
P3 null: <div></div>
P3 false: <div></div>
P3 function: <div><em>Hello, Ada</em></div>
\end{lstlisting}

The function case works because \passthrough{\lstinline!ChildShowcase!}
checks whether \passthrough{\lstinline!children!} is a function and
invokes it. Passing a function as ordinary children without invoking it
is not how React renders content.

\subsubsection{P4: separate children from the other
props}\label{p4-separate-children-from-the-other-props}

The pattern \passthrough{\lstinline!(\{ title, children, ...rest \})!}
extracts the values the component needs and collects the remaining
properties into \passthrough{\lstinline!rest!}. Spreading
\passthrough{\lstinline!rest!} onto \passthrough{\lstinline!<section>!}
forwards those values to the DOM element.

\textbf{Rendered output:}

\begin{lstlisting}
P4 forwarded: <section id="agenda" class="panel"><h2>Agenda</h2>Items</section>
P4 rest wins: <section id="from-rest"><h2>Agenda</h2>Items</section>
P4 explicit wins: <section id="fixed-id"><h2>Agenda</h2>Items</section>
\end{lstlisting}

In \passthrough{\lstinline!<section id="fixed-id" \{...rest\}>!}, a
later \passthrough{\lstinline!id!} from \passthrough{\lstinline!rest!}
overwrites the earlier explicit value. In
\passthrough{\lstinline!<section \{...rest\} id="fixed-id">!}, the later
explicit value wins. The order of spread and explicit props therefore
matters.

\subsubsection{P5: named header/footer props or
children?}\label{p5-named-headerfooter-props-or-children}

A component with \passthrough{\lstinline!header!} and
\passthrough{\lstinline!footer!} props communicates that it owns those
two slots and keeps the middle structure fixed. A component that accepts
\passthrough{\lstinline!children!} lets the caller compose the whole
internal layout. Choose named props when the slots are fixed and clearly
named; choose children when flexible composition is the point.

\textbf{Rendered output:}

\begin{lstlisting}
P5 named: <article><header>Course</header><main>Shared content</main><footer>2026</footer></article>
P5 children: <article><header>Course</header><main>Shared content</main><footer>2026</footer></article>
\end{lstlisting}

\subsubsection{P6: defaults and null}\label{p6-defaults-and-null}

A parameter default such as
\passthrough{\lstinline!\{ name = 'Guest' \}!} applies when
\passthrough{\lstinline!name!} is \passthrough{\lstinline!undefined!}
(including when the prop is omitted). It does not apply when
\passthrough{\lstinline!name!} is \passthrough{\lstinline!null!}. The
example uses \passthrough{\lstinline!name ?? 'nobody'!} in the rendered
expression to make the explicit \passthrough{\lstinline!null!} case
visible.

\textbf{Rendered output:}

\begin{lstlisting}
P6 missing: <p>Hello, Guest!</p>
P6 undefined: <p>Hello, Guest!</p>
P6 null: <p>Hello, nobody!</p>
\end{lstlisting}

\subsubsection{P7: typing props and
children}\label{p7-typing-props-and-children}

The JavaScript runtime counterpart in
\passthrough{\lstinline!js/props.jsx!} forwards
\passthrough{\lstinline!type!} and \passthrough{\lstinline!aria-label!}
to a native button, but JavaScript itself does not statically reject an
unknown or missing prop. Three common TypeScript choices are shown in
\passthrough{\lstinline!ts/props.tsx!}:

\begin{itemize}
\tightlist
\item
  \passthrough{\lstinline!children?: ReactNode!} explicitly adds an
  optional children property to a custom props type. This allows text,
  numbers, elements, arrays and the other renderable React node values.
\item
  \passthrough{\lstinline!PropsWithChildren<\{ disabled?: boolean \}>!}
  adds an optional \passthrough{\lstinline!children!} property to the
  supplied props shape. It is convenient when a component accepts normal
  React children.
\item
  \passthrough{\lstinline!ComponentProps<'button'>!} reuses the props
  type of the native button element. It includes button attributes such
  as \passthrough{\lstinline!type!}, \passthrough{\lstinline!disabled!},
  and \passthrough{\lstinline!aria-label!}, as well as children.
\end{itemize}

\textbf{Rendered output:}

\begin{lstlisting}
P7 JavaScript: <button type="submit" aria-label="Send">Send</button>
P7 manual: <button disabled="">Save</button>
P7 PropsWithChildren: <button>Cancel</button>
P7 ComponentProps: <button type="submit" aria-label="Send">Send</button>
\end{lstlisting}

If a custom props type does not declare
\passthrough{\lstinline!children!}, TypeScript rejects JSX content for
that component. Add \passthrough{\lstinline!children?: ReactNode!} or
use \passthrough{\lstinline!PropsWithChildren!} when children are
intended. Do not use \passthrough{\lstinline!any!} to silence a missing
prop type.

\subsection{TypeScript topics T1-T5}\label{typescript-topics-t1-t5}

These examples are in \passthrough{\lstinline!ts/types.ts!}.

\begin{itemize}
\tightlist
\item
  \textbf{T1 --- rename is not a type annotation.} In
  \passthrough{\lstinline!const \{ x: horizontal \} = point!},
  \passthrough{\lstinline!x!} is the source property and
  \passthrough{\lstinline!horizontal!} is the new variable. To type the
  pattern, put the annotation after the closing brace:
  \passthrough{\lstinline!const \{ y: vertical \}: \{ y: number \} = point!}.
\item
  \textbf{T2 --- tuple return and \passthrough{\lstinline!as const!}.} A
  return type such as \passthrough{\lstinline![number, string]!} says
  the first position is a number and the second is a string.
  \passthrough{\lstinline!as const!} makes an array literal a readonly
  tuple with literal types, rather than a general array whose elements
  may have a union type.
\item
  \textbf{T3 --- \passthrough{\lstinline!noUncheckedIndexedAccess!}.}
  With this compiler option, an array item may be missing. Therefore
  \passthrough{\lstinline!const [firstItem] = items!} has type
  \passthrough{\lstinline!string | undefined!} when
  \passthrough{\lstinline!items!} is \passthrough{\lstinline!string[]!}.
  Use a check or a fallback such as
  \passthrough{\lstinline!firstItem ?? 'fallback item'!}.
\item
  \textbf{T4 --- optional fields and defaults.} An optional
  \passthrough{\lstinline!nickname?: string | null!} can be
  \passthrough{\lstinline!string!}, \passthrough{\lstinline!null!}, or
  \passthrough{\lstinline!undefined!}. Destructuring without a default
  preserves all three possibilities. A default removes
  \passthrough{\lstinline!undefined!} from the local variable type, but
  it does not remove \passthrough{\lstinline!null!}.
\item
  \textbf{T5 --- a type error and its fix.}
  \passthrough{\lstinline!square!} requires a number, so
  \passthrough{\lstinline!square('4')!} produces TS2345:
  \passthrough{\lstinline!Argument of type 'string' is not assignable to parameter of type 'number'.!}
  The example marks this one intentional line with
  \passthrough{\lstinline!@ts-expect-error!}; the corrected call is
  \passthrough{\lstinline!square(4)!}.
\end{itemize}

\subsection{JavaScript versus
TypeScript}\label{javascript-versus-typescript}

\begin{longtable}[]{@{}llll@{}}
\toprule\noalign{}
Form/topic & What stays the same & What TypeScript adds & What can go
wrong \\
\midrule\noalign{}
\endhead
\bottomrule\noalign{}
\endlastfoot
Array positions (1) & \passthrough{\lstinline![a, b]!} reads positions 0
and 1 & Tuple return type such as
\passthrough{\lstinline![number, number]!} & Wrong tuple position type
or possible \passthrough{\lstinline!undefined!} \\
Skip and rest (2) & \passthrough{\lstinline![, second, ...rest]!} &
Tuple and rest element types & Incorrect tuple shape or a rest type that
is too narrow \\
Swap (3) & \passthrough{\lstinline![a, b] = [b, a]!} & Variables must
have compatible declared types & Incompatible assignment types \\
Array defaults (4) & Defaults apply to
\passthrough{\lstinline!undefined!} & Optional tuple positions and
inferred fallback types & A value can still be
\passthrough{\lstinline!undefined!} if no fallback is provided \\
Object properties (5) & \passthrough{\lstinline!\{ title \}!} reads the
\passthrough{\lstinline!title!} property & An object type must contain a
compatible property & Missing property or incompatible property type \\
Rename (6) & \passthrough{\lstinline!\{ title: heading \}!} creates
\passthrough{\lstinline!heading!} & The source property and new variable
receive inferred/declared types & Mistaking
\passthrough{\lstinline!title!} for the new variable; invalid source
key \\
Object defaults (7) &
\passthrough{\lstinline!\{ location = 'Library' \}!} uses a fallback for
\passthrough{\lstinline!undefined!} & Optional property types and a
narrowed local type & \passthrough{\lstinline!null!} remains possible if
included in the type \\
Object rest (8) & \passthrough{\lstinline!\{ id, ...rest \}!} collects
other properties & Rest receives the remaining inferred object type &
Using a property that was removed by the pattern \\
Nested patterns (9) & Nested object/array shape is read & Nested
property types are checked & Possibly missing nested property or
\passthrough{\lstinline!undefined!} array item \\
Computed key (10) & \passthrough{\lstinline!\{ [key]: value \}!}
evaluates a property key & Key type can be narrowed to known keys & A
broad string key may not identify a known property \\
Parameters (11) & A parameter pattern extracts options & Type annotation
follows the whole pattern & Wrong annotation placement or missing
required argument fields \\
Loops (12) & Entry pairs can be destructured &
\passthrough{\lstinline!Object.entries!} and
\passthrough{\lstinline!Map<K, V>!} provide key/value types & Values may
be wider than expected; unchecked array indexing remains possible \\
\end{longtable}

\subsection{Output and verification}\label{output-and-verification}

The JavaScript-only files were run during development and produced the
outputs shown in the field guide and pitfalls section. The full
\passthrough{\lstinline!output.txt!} must be generated on your own
machine after dependencies are installed and after
\passthrough{\lstinline!npm run all!} succeeds:

\begin{lstlisting}[language=bash]
npm install
npm run all
npm run check
npm run all > output.txt
\end{lstlisting}

Do not treat a partial run as the final
\passthrough{\lstinline!output.txt!}. Confirm that the React output and
TypeScript examples also appear in the completed file.

\subsection{My own log}\label{my-own-log}

These are the issues and surprises I actually encountered while setting
up and running this project on my Mac.

\begin{enumerate}
\def\labelenumi{\arabic{enumi}.}
\tightlist
\item
  \textbf{The Node.js version selected by my Terminal was too old.}
  Exact output: \passthrough{\lstinline!node -v!} returned
  \passthrough{\lstinline!v20.17.0!}, while the project instructions
  require Node.js 20.19 or newer. I initially expected the installed
  Node version to be used automatically. I checked
  \passthrough{\lstinline!which node!} and found
  \passthrough{\lstinline!/opt/anaconda3/bin/node!}. I learned that the
  order of folders in \passthrough{\lstinline!PATH!} determines which
  executable the Terminal finds first.
\item
  \textbf{Anaconda was taking priority over the newer Node
  installation.} Exact output: \passthrough{\lstinline!which node!}
  returned \passthrough{\lstinline!/opt/anaconda3/bin/node!}; running
  \passthrough{\lstinline!/usr/local/bin/node -v!} returned
  \passthrough{\lstinline!v26.13.1!}. I ran
  \passthrough{\lstinline!export PATH="/usr/local/bin:$PATH"!} and
  checked again; \passthrough{\lstinline!node -v!} then returned
  \passthrough{\lstinline!v26.13.1!}. I learned how to temporarily
  change \passthrough{\lstinline!PATH!} for the current Terminal session
  and why \passthrough{\lstinline!which node!} is useful for diagnosing
  version problems.
\item
  \textbf{I entered a Terminal command while I was in the Node.js REPL.}
  Exact error:
  \passthrough{\lstinline!ReferenceError: ls is not defined!}. I
  expected \passthrough{\lstinline!ls!} to list the files, but the Node
  REPL interpreted it as JavaScript. I exited the REPL and ran the
  command at the normal shell prompt. I learned to distinguish the Node
  prompt from the Terminal shell prompt.
\item
  \textbf{I tried to use shell syntax inside the Node.js REPL.} Exact
  error:
  \passthrough{\lstinline!SyntaxError: Unexpected identifier '$PATH'!}.
  I expected \passthrough{\lstinline!$PATH!} to be understood as a shell
  variable, but the REPL parses JavaScript. I moved the
  \passthrough{\lstinline!export PATH=...!} command to the normal
  Terminal prompt. I learned that shell commands and JavaScript
  expressions belong in different environments.
\item
  \textbf{\passthrough{\lstinline!npm install!} displayed install-script
  warnings, but the installation completed.} Exact warning named
  \passthrough{\lstinline!esbuild@0.28.2!} and
  \passthrough{\lstinline!fsevents@2.3.3!} as packages with install
  scripts not yet covered by \passthrough{\lstinline!allowScripts!}; npm
  still reported that packages were added and
  \passthrough{\lstinline!found 0 vulnerabilities!}. I initially was
  unsure whether this meant installation had failed. I continued by
  running the project and confirmed that
  \passthrough{\lstinline!npm run all!} worked. I learned to distinguish
  a warning from a failed command and to verify by running the required
  scripts.
\item
  \textbf{React renders some different JavaScript child values as empty
  output.} Actual outputs:
  \passthrough{\lstinline!P3 undefined: <div></div>!},
  \passthrough{\lstinline!P3 null: <div></div>!}, and
  \passthrough{\lstinline!P3 false: <div></div>!}. I expected these
  values might appear in the markup, but
  \passthrough{\lstinline!renderToStaticMarkup!} produced an empty
  \passthrough{\lstinline!<div>!}. I learned that these values are valid
  React children but do not produce visible text or elements in the
  rendered HTML.
\item
  \textbf{A destructuring default handles
  \passthrough{\lstinline!undefined!}, but not
  \passthrough{\lstinline!null!}.} Actual outputs:
  \passthrough{\lstinline"P6 undefined: <p>Hello, Guest!</p>"} and
  \passthrough{\lstinline"P6 null: <p>Hello, nobody!</p>"}. I expected
  the default value to cover both cases. The output showed that an
  omitted or \passthrough{\lstinline!undefined!} property activates the
  default, while explicit \passthrough{\lstinline!null!} remains a
  separate value. I learned to handle \passthrough{\lstinline!null!}
  explicitly when the application needs a fallback for it.
\end{enumerate}

\subsection{Bonus 2: Depth}\label{bonus-2-depth}

The depth experiments are implemented in
\passthrough{\lstinline!js/depth.mjs!} and
\passthrough{\lstinline!ts/depth.ts!}. They extend the basic forms with
iterable protocols, evaluation order, assignment targets, and TypeScript
type-level patterns. Run them together with the main examples using
\passthrough{\lstinline!npm run all!}.

\subsubsection{B2.1 Strings are
iterable}\label{b21-strings-are-iterable}

JavaScript strings are iterable, so an array pattern can consume their
characters in iteration order. In this example, the first two characters
are assigned to \passthrough{\lstinline!first!} and
\passthrough{\lstinline!second!}, while the rest pattern collects the
remaining characters into an array.

Actual output:

\begin{lstlisting}
first two: A d
remaining: a| |L|o|v|e|l|a|c|e
\end{lstlisting}

\subsubsection{B2.2 Set and Map are
iterable}\label{b22-set-and-map-are-iterable}

A \passthrough{\lstinline!Set!} yields its values in insertion order. A
\passthrough{\lstinline!Map!} yields entry pairs, where each entry is a
two-element array containing a key and its value. This makes nested
destructuring useful: \passthrough{\lstinline![key, value]!} extracts
each pair\textquotesingle s two parts.

Actual output:

\begin{lstlisting}
Set: teal coral, navy
Map first entry: A1 Mina
Map remaining entries: [["A2","Eren"]]
\end{lstlisting}

\subsubsection{B2.3 Destructuring consumes an
iterator}\label{b23-destructuring-consumes-an-iterator}

Array destructuring reads from an iterator; it does not merely inspect
an existing array. A generator is a one-use iterator. The pattern
requests two values, and because the pattern finishes before the
generator is naturally exhausted, iterator closing occurs. The generator
is then closed, so a later \passthrough{\lstinline!next()!} does not
produce another task.

Actual output:

\begin{lstlisting}
destructured: T-01 T-02
next value after destructuring: undefined
\end{lstlisting}

\subsubsection{\texorpdfstring{B2.4 Defaults run only for
\texttt{undefined}, from left to
right}{B2.4 Defaults run only for undefined, from left to right}}\label{b24-defaults-run-only-for-undefined-from-left-to-right}

A destructuring default initializer is used only when the incoming value
is \passthrough{\lstinline!undefined!}. It does not replace
\passthrough{\lstinline!null!}. Initializers are evaluated when needed
and in left-to-right order. A later binding can refer to an earlier
binding, but referring to a later \passthrough{\lstinline!let!} binding
before it has initialized causes a temporal-dead-zone
\passthrough{\lstinline!ReferenceError!}.

Actual output:

\begin{lstlisting}
values: Read todo
fallback calls: 0
null keeps its value: null
fallback calls after null: 0
later default can use earlier binding: 10 20
reversed order: ReferenceError (temporal dead zone)
\end{lstlisting}

The zero fallback-call count demonstrates that the fallback function was
not evaluated when a value was already present. The
\passthrough{\lstinline!null!} example demonstrates why
\passthrough{\lstinline!??!} or an explicit null check may be needed
when both \passthrough{\lstinline!null!} and
\passthrough{\lstinline!undefined!} should trigger a fallback.

\subsubsection{B2.5 Destructuring assignment can target object
properties}\label{b25-destructuring-assignment-can-target-object-properties}

A destructuring assignment pattern can assign directly to existing
object properties rather than declaring new variables. When an object
pattern is used as an assignment statement, wrap the pattern in
parentheses so JavaScript does not parse the opening brace as a block.

Actual output:

\begin{lstlisting}
updated point: 3 7
\end{lstlisting}

\subsubsection{\texorpdfstring{B2.6 \texttt{for...of} can destructure
entry pairs; \texttt{for...in} does
not}{B2.6 for...of can destructure entry pairs; for...in does not}}\label{b26-forof-can-destructure-entry-pairs-forin-does-not}

\passthrough{\lstinline!Object.entries(object)!} returns an array of
\passthrough{\lstinline![key, value]!} pairs.
\passthrough{\lstinline!for...of!} iterates those pairs, so
\passthrough{\lstinline!for (const [key, value] of Object.entries(object))!}
can extract both values. By contrast, \passthrough{\lstinline!for...in!}
iterates enumerable property keys; it does not automatically produce
\passthrough{\lstinline![key, value]!} pairs.

Actual output:

\begin{lstlisting}
title=Read
status=todo
for...in yields property names, not [key, value] pairs:
title
status
\end{lstlisting}

\subsubsection{TypeScript type-level
destructuring}\label{typescript-type-level-destructuring}

TypeScript conditional types can use \passthrough{\lstinline!infer!} to
extract parts of tuple types. \passthrough{\lstinline!First<T>!} infers
the first tuple element, while \passthrough{\lstinline!Rest<T>!} infers
the remaining elements. The built-in
\passthrough{\lstinline!Parameters<F>!} utility type extracts a
function\textquotesingle s parameter types as a tuple.

Actual output:

\begin{lstlisting}
type-level First: 41
type-level Rest: [ 'MIS', true ]
Parameters<F>: [ 'T-01', false ]
\end{lstlisting}

\subsubsection{Verification}\label{verification}

On my Mac, \passthrough{\lstinline!npm run all!} completed and printed
all six B2 JavaScript experiments and the TypeScript type-level results.
I then ran \passthrough{\lstinline!npm run all > output.txt!} to save
the complete output and \passthrough{\lstinline!npm run check!}; the
TypeScript compiler completed without reporting errors. I also used
\passthrough{\lstinline!grep!} to confirm that
\passthrough{\lstinline!output.txt!} contains the B2.1-B2.6 headings,
\passthrough{\lstinline!type-level First!},
\passthrough{\lstinline!type-level Rest!}, and
\passthrough{\lstinline!Parameters<F>!}.

\subsection{Bonus 1: LaTeX}\label{bonus-1-latex}

This complete field guide is also provided as
\passthrough{\lstinline!destructuring-field-guide.tex!} and the compiled
\passthrough{\lstinline!destructuring-field-guide.pdf!}. The LaTeX
source was generated from this full Markdown guide with Pandoc,
including the field-guide entries, pitfalls, React props and children
discussion, JavaScript-versus-TypeScript comparison table, and the Bonus
2 depth experiments. The PDF includes a table of contents and code
listings.

Commands used to generate and compile the LaTeX source:

\begin{lstlisting}[language=bash]
pandoc README.md --from=gfm --to=latex --standalone --toc --number-sections --listings -V geometry:margin=1in -o destructuring-field-guide.tex
xelatex -interaction=nonstopmode -halt-on-error destructuring-field-guide.tex
xelatex -interaction=nonstopmode -halt-on-error destructuring-field-guide.tex
\end{lstlisting}

The second XeLaTeX run updates the table of contents and
cross-references. The \passthrough{\lstinline!.tex!} source and compiled
PDF are submitted together so the document can be inspected and rebuilt
from source.

\subsection{Bonus 3: Real code}\label{bonus-3-real-code}

This section examines destructuring in real projects rather than only in
isolated teaching examples. Each case identifies the source location,
the destructuring form, an equivalent version without destructuring, and
why the pattern is useful or potentially less readable.

\subsubsection{\texorpdfstring{R1 --- Task Manager: \texttt{useReducer}
result (array
destructuring)}{R1 --- Task Manager: useReducer result (array destructuring)}}\label{r1--task-manager-usereducer-result-array-destructuring}

\textbf{Source:} Assignment 3, \passthrough{\lstinline!src/App.tsx!},
line 85.

\begin{lstlisting}
const [state, dispatch] = useReducer(taskReducer, initialState)
\end{lstlisting}

\textbf{Form:} Form 1 --- array values by position.

\passthrough{\lstinline!useReducer!} returns a pair: the current state
and the dispatch function. The two array positions have a stable,
documented meaning, so positional destructuring gives each value a
useful local name.

\textbf{Equivalent without destructuring:}

\begin{lstlisting}
const reducerResult = useReducer(taskReducer, initialState)
const state = reducerResult[0]
const dispatch = reducerResult[1]
\end{lstlisting}

\subsubsection{R2 --- Task Manager: destructuring component
props}\label{r2--task-manager-destructuring-component-props}

\textbf{Source:} Assignment 3,
\passthrough{\lstinline!src/TaskRow.tsx!}, line 15.

\begin{lstlisting}
export default function TaskRow({ task, onNext, onDelete }: TaskRowProps) {
\end{lstlisting}

\textbf{Form:} Form 5 + Form 11 --- object properties by name,
destructured in a function parameter.

React passes the component one props object. This parameter pattern
extracts the three named properties directly. In this component, the
declared props are \passthrough{\lstinline!task!},
\passthrough{\lstinline!onNext!}, and
\passthrough{\lstinline!onDelete!}; the component does not declare or
receive a \passthrough{\lstinline!children!} prop.

\textbf{Equivalent without destructuring:}

\begin{lstlisting}
export default function TaskRow(props: TaskRowProps) {
  const task = props.task
  const onNext = props.onNext
  const onDelete = props.onDelete
\end{lstlisting}

\subsubsection{\texorpdfstring{R3 --- Task Manager: extracting
\texttt{dispatch} from
props}{R3 --- Task Manager: extracting dispatch from props}}\label{r3--task-manager-extracting-dispatch-from-props}

\textbf{Source:} Assignment 3, \passthrough{\lstinline!src/App.tsx!},
line 17.

\begin{lstlisting}
function AddForm({ dispatch }: AddFormProps) {
\end{lstlisting}

\textbf{Form:} Form 5 + Form 11 --- object property destructuring in a
parameter.

Only one property is needed by \passthrough{\lstinline!AddForm!}, so
destructuring avoids repeating \passthrough{\lstinline!props.dispatch!}
throughout the component.

\textbf{Equivalent without destructuring:}

\begin{lstlisting}
function AddForm(props: AddFormProps) {
  const dispatch = props.dispatch
\end{lstlisting}

\subsubsection{R4 --- TypeScript compiler: renaming object
properties}\label{r4--typescript-compiler-renaming-object-properties}

\textbf{Source:} TypeScript 5.8.3 package,
\passthrough{\lstinline!node\_modules/typescript/lib/typescript.js!}.
The corresponding readable upstream source is
\passthrough{\lstinline!src/compiler/program.ts!}, lines 3452-3454, in
the TypeScript \passthrough{\lstinline!v5.8.3!} tag:
\textless{}\url{https://github.com/microsoft/TypeScript/blob/v5.8.3/src/compiler/program.ts\#L3452-L3454}.

\begin{lstlisting}
const { line: firstLine, character: firstLineChar } =
  getLineAndCharacterOfPosition(file, start)
const { line: lastLine, character: lastLineChar } =
  getLineAndCharacterOfPosition(file, start + length)
\end{lstlisting}

\textbf{Form:} Form 6 --- object property renaming.

The function returns an object with \passthrough{\lstinline!line!} and
\passthrough{\lstinline!character!} properties. The code renames them to
\passthrough{\lstinline!firstLine!},
\passthrough{\lstinline!firstLineChar!},
\passthrough{\lstinline!lastLine!}, and
\passthrough{\lstinline!lastLineChar!}, which makes their roles clear in
the surrounding formatting logic.

\textbf{Equivalent without destructuring:}

\begin{lstlisting}
const firstPosition = getLineAndCharacterOfPosition(file, start)
const firstLine = firstPosition.line
const firstLineChar = firstPosition.character
const lastPosition = getLineAndCharacterOfPosition(file, start + length)
const lastLine = lastPosition.line
const lastLineChar = lastPosition.character
\end{lstlisting}

To record the exact line in the installed bundled file on this machine,
run the following command from the Assignment 4 folder after
\passthrough{\lstinline!npm install!}:

\begin{lstlisting}[language=bash]
grep -n -F 'getLineAndCharacterOfPosition(file, start)' node_modules/typescript/lib/typescript.js | head
\end{lstlisting}

The upstream source line numbers above are pinned to TypeScript v5.8.3;
the generated bundle may have different line numbering.

\subsubsection{\texorpdfstring{R5 --- React Navigation: a React
component receives \texttt{children} as part of
props}{R5 --- React Navigation: a React component receives children as part of props}}\label{r5--react-navigation-a-react-component-receives-children-as-part-of-props}

\textbf{Source:} React Navigation,
\passthrough{\lstinline!packages/stack/src/views/Stack/CardContainer.tsx!},
lines 931 and 935-998:
\textless{}\url{https://github.com/react-navigation/react-navigation/blob/main/packages/stack/src/views/Stack/CardContainer.tsx\#L931-L998}.

The \passthrough{\lstinline!Props!} type declares
\passthrough{\lstinline!children: React.ReactNode!} at line 931.
\passthrough{\lstinline!CardContainerInner!} then destructures
\passthrough{\lstinline!children!} alongside the other properties of the
same props object:

\begin{lstlisting}
function CardContainerInner({
  interpolationIndex,
  index,
  active,
  // ...other props
  scene,
  children,
}: Props) {
\end{lstlisting}

\textbf{Form:} Form 5 + Form 11 --- object properties destructured in a
function parameter.

This is a concrete example of the distinction between ordinary props and
\passthrough{\lstinline!children!}: \passthrough{\lstinline!children!}
is itself a property on the props object. In JSX, nested content is
supplied through that property; it is not passed as a second function
argument.

\textbf{Equivalent without destructuring:}

\begin{lstlisting}
function CardContainerInner(props: Props) {
  const interpolationIndex = props.interpolationIndex
  const index = props.index
  const active = props.active
  // Read other values as props.someProperty when needed.
  const scene = props.scene
  const children = props.children
\end{lstlisting}

This component has a large props shape, but that does not mean every
component should destructure every property. It is useful to extract the
values used frequently while keeping the relationship to
\passthrough{\lstinline!props!} explicit when that improves clarity.

\subsubsection{R6 --- React Navigation: when destructuring many options
can make code harder to
scan}\label{r6--react-navigation-when-destructuring-many-options-can-make-code-harder-to-scan}

\textbf{Source:} React Navigation,
\passthrough{\lstinline!packages/stack/src/views/Stack/CardContainer.tsx!},
lines 1105-1135:
\textless{}\url{https://github.com/react-navigation/react-navigation/blob/main/packages/stack/src/views/Stack/CardContainer.tsx\#L1105-L1135}.

\begin{lstlisting}
const {
  presentation,
  animation,
  cardOverlay,
  cardOverlayEnabled,
  cardShadowEnabled,
  cardStyle,
  cardStyleInterpolator,
  gestureDirection,
  gestureEnabled,
  gestureResponseDistance,
  gestureVelocityImpact,
  headerMode,
  headerShown,
  transitionSpec,
} = scene.descriptor.options
\end{lstlisting}

\textbf{Form:} Form 5 --- object property destructuring.

\textbf{Why this can be worse:} extracting a long list of similarly
named options creates many local variables and makes it harder to see
that they all come from
\passthrough{\lstinline!scene.descriptor.options!}. If only a few values
are needed in a section, or if the source object is important context,
the long list can add visual noise. This is a trade-off, not a rule that
destructuring is always bad.

\textbf{Equivalent without destructuring:}

\begin{lstlisting}
const options = scene.descriptor.options
// Use options.presentation, options.animation, options.cardStyle, etc.
\end{lstlisting}

Keeping the object grouped can make the origin of each setting more
obvious. Destructuring is most helpful when the extracted names improve
local readability; it is not automatically better just because it is
shorter.

\subsubsection{Real-code comparison
summary}\label{real-code-comparison-summary}

\begin{longtable}[]{@{}llll@{}}
\toprule\noalign{}
Case & Source & Form & Main lesson \\
\midrule\noalign{}
\endhead
\bottomrule\noalign{}
\endlastfoot
R1 & Assignment 3 \passthrough{\lstinline!src/App.tsx:85!} & Form 1 & A
tuple-like return value has meaningful positions. \\
R2 & Assignment 3 \passthrough{\lstinline!src/TaskRow.tsx:15!} & Forms 5
and 11 & A React component receives one props object. \\
R3 & Assignment 3 \passthrough{\lstinline!src/App.tsx:17!} & Forms 5 and
11 & Extract only the prop a component needs. \\
R4 & TypeScript 5.8.3 \passthrough{\lstinline!program.ts:3452-3454!} /
installed package bundle & Form 6 & Renaming gives generic property
names a local meaning. \\
R5 & React Navigation
\passthrough{\lstinline!CardContainer.tsx:931, 935-998!} & Forms 5 and
11 & \passthrough{\lstinline!children!} is a property of the props
object. \\
R6 & React Navigation
\passthrough{\lstinline!CardContainer.tsx:1105-1135!} & Form 5 & A long
destructuring list can reduce readability. \\
\end{longtable}

The Task Manager examples refer to the Assignment 3 source files. The
TypeScript and React Navigation examples link to their public source
repositories so the code can be checked independently.

\end{document}
